\chapter{Convolutional Neural Networks}

\section{Introduction}

Convolutional neural networks (CNNs) are designed to exploit spatial structure in data, particularly images. Unlike fully connected networks, CNNs incorporate locality, weight sharing, and translation equivariance directly into their architecture.

This chapter develops CNNs from a mathematical perspective. We interpret convolution as a structured linear operator, formalize translation equivariance, and discuss pooling and receptive fields. We also briefly connect convolution to the Fourier domain.

\section{Convolution as a Local Linear Operator}

A convolution layer applies a linear operator with a special structure.

\subsection{Discrete Convolution}

For a 1D signal $x \in \mathbb{R}^n$ and kernel $k \in \mathbb{R}^r$, the discrete convolution is:
\[
(y * k)_i = \sum_{j=0}^{r-1} k_j \, x_{i-j}.
\]

In 2D (for images), with input $x$ and kernel $k$:
\[
(y * k)_{i,j} = \sum_{u,v} k_{u,v} \, x_{i-u, j-v}.
\]

\subsection{Structured Linear Map}

Convolution can be written as a matrix-vector product:
\[
y = K x,
\]
where $K$ is a Toeplitz (or block Toeplitz) matrix. Unlike a fully connected layer, $K$ has:
\begin{itemize}
    \item repeated entries (weight sharing),
    \item sparse local structure.
\end{itemize}

Thus convolution is a linear map with strong structural constraints.

\section{Weight Sharing and Equivariance}

\subsection{Weight Sharing}

In convolution, the same kernel is applied across all spatial locations. This drastically reduces the number of parameters compared to fully connected layers.

\subsection{Translation Equivariance}

\begin{definition}[Translation Operator]
Let $T_\tau$ shift an input by $\tau$. For example:
\[
(T_\tau x)_i = x_{i-\tau}.
\]
\end{definition}

\begin{proposition}[Equivariance of Convolution]
For a convolution operator $C$, we have:
\[
C(T_\tau x) = T_\tau (C x).
\]
\end{proposition}

\subsection{Interpretation}

A shift in the input produces a corresponding shift in the output. Thus convolution preserves spatial structure.

This is a key inductive bias:
\begin{itemize}
    \item features detected in one location are detected everywhere,
    \item the model is robust to translations.
\end{itemize}

\section{Pooling, Receptive Fields, and Hierarchy}

\subsection{Pooling}

Pooling reduces spatial resolution:
\begin{itemize}
    \item max pooling selects the maximum in a local region,
    \item average pooling computes local averages.
\end{itemize}

Pooling introduces:
\begin{itemize}
    \item invariance to small translations,
    \item reduced computational cost.
\end{itemize}

\subsection{Receptive Fields}

The receptive field of a unit is the region of the input that influences it.

\begin{itemize}
    \item early layers: small receptive fields (local features),
    \item deeper layers: larger receptive fields (global structure).
\end{itemize}

Stacking convolution and pooling layers increases the effective receptive field.

\subsection{Hierarchical Representation}

CNNs build hierarchical features:
\begin{itemize}
    \item edges and textures in early layers,
    \item shapes and parts in intermediate layers,
    \item objects in deeper layers.
\end{itemize}

This mirrors the compositional structure of visual data.

\section{CNN Architectures in Vision}

Typical CNN architectures combine:
\begin{itemize}
    \item convolution layers,
    \item nonlinear activations,
    \item pooling or strided convolutions,
    \item normalization,
    \item fully connected or global pooling layers.
\end{itemize}

Modern architectures often incorporate:
\begin{itemize}
    \item residual connections,
    \item bottleneck blocks,
    \item deeper hierarchies.
\end{itemize}

These design choices reflect tradeoffs between expressivity, efficiency, and trainability.

\section{Spectral/Fourier Perspective}

Convolution has a natural interpretation in the frequency domain.

\subsection{Convolution Theorem}

For suitable signals:
\[
\mathcal{F}(x * k) = \mathcal{F}(x) \cdot \mathcal{F}(k),
\]
where $\mathcal{F}$ denotes the Fourier transform.

\subsection{Interpretation}

\begin{itemize}
    \item convolution in space corresponds to multiplication in frequency,
    \item kernels act as frequency filters,
    \item CNNs can be viewed as learning filters in the spectral domain.
\end{itemize}

For example:
\begin{itemize}
    \item edge detectors emphasize high frequencies,
    \item smoothing filters suppress high frequencies.
\end{itemize}

Thus CNNs can be interpreted as adaptive filter banks.

\section{Worked Examples}

\subsection{Feature Map Size}

For input size $n$, kernel size $r$, and stride $s$, the output size is approximately:
\[
\left\lfloor \frac{n - r}{s} \right\rfloor + 1.
\]

Padding can adjust this to preserve spatial dimensions.

\subsection{Equivariance Example}

Let $x$ be a signal and $T_\tau$ a shift. Then:
\[
(Cx)_i = \sum_j k_j x_{i-j},
\]
\[
(C(T_\tau x))_i = \sum_j k_j x_{i-j-\tau} = (Cx)_{i-\tau}.
\]
Thus:
\[
C(T_\tau x) = T_\tau(Cx).
\]

\section{A Unifying Perspective}

CNNs impose strong structural assumptions:
\begin{itemize}
    \item locality: interactions are local,
    \item weight sharing: same features across space,
    \item equivariance: structure preserved under translation.
\end{itemize}

These assumptions reduce the hypothesis space while aligning it with natural image structure.

From a linear-algebra viewpoint, convolution is a structured matrix. From a signal-processing viewpoint, it is filtering. From a learning viewpoint, it encodes inductive bias.

\section{Summary}

In this chapter we developed convolutional neural networks from a mathematical perspective. We defined discrete convolution as a structured linear operator and showed how weight sharing leads to translation equivariance. We discussed pooling, receptive fields, and hierarchical feature construction.

We also introduced the Fourier viewpoint, interpreting convolution as frequency-domain filtering. CNNs illustrate how architectural structure can encode powerful inductive biases for specific data domains.